\section*{Alban : de l'Édit de Nantes à la Guerre de Rohan}

La paix étant relativement retrouvée, on s'organise. Monsieur Legoux de la Berchère fait de Notre Dame d'Alban un chef lieu de « claverie » et de district (circonscription administrative et pastorale). Nous disons actuellement un « doyenné ». Le curé de N.D. d'Alban prend le titre de « Vicaire forain ». Il l'a encore de nos jours et remplace le titre d'archiprêtre, jadis porté par le curé de Négremont. Le titre est-il attaché nécessairement au lieu d'Alban ? On a pu le croire jusqu'à une date récente. Au point que cela causa quelque étonnement en 1965 quand le curé de Trébas en porta le titre, le doyen d'Alban ayant offert sa démission de doyen. Mais l'histoire apprend qu'il n'y a rien de nouveau sous le soleil. Ainsi, deux siècles plus tôt, en 1765, l'abbé Jean-Pierre Rouziès, curé de Saint Projet de Paulin, est « vicaire forain du district d'Alban », si du moins nous en croyons l'acte officiel du 4 août 1765 consigné dans le registre de Saint Jean de Janes. Ce canton ecclésiastique ou « chambre d'Alban » est arrenté\footnote{???} à un neveu de l'Archevêque pour la somme de six cents écus et quatre sols par an. Qui fit une bonne affaire, Monsieur Delbène ou son parent ?... Problème.

Notre cité est en piteux état et s'essaie à relever ses ruines. Le cadastre de 1598 y mentionne encore une porte basse vers le chemin d'Alban à Albi. Dans le sens opposé se dresse la porte haute. On peut penser que c'est pour ce motif que se désignait, encore il y a vingt ans, la partie est du village, sous le vocable de « Cap d'Alban ». Trente ans en arrière, les anciens peuvent se souvenir que « lou cap d'Alban » avait le privilège d'organiser une cavalcade à l'occasion des fêtes du village.